\documentclass{article}
\title{論理の組み立て方 — 仕事の主張を崩されない形にする}
\author{TeX 図書館}

\begin{document}

\section{主張を1つに絞る — 何を言ってほしいのか}

仕事の文章が崩される理由は、論理の細部より先に、\textbf{出発点の曖昧さ}にあります。何を言ってほしいのかが定まっていないと、相手は好きな箇所だけを潰して終わります。

報告も提案もメールも会議も、構造は同じです。主張が1つ、根拠がそろっている、想定反論に答えている。この3つが揃えば、主張は崩れにくくなります。まず主張から始めます。

\begin{quote}
\textbf{【例】}\ ある報告メールの冒頭。「今週の問い合わせ状況についてです。あと人員の話もしたいです。あと来期の予算のことも相談させてください。とりあえず、状況から報告します。」
\end{quote}

このメールには主張がありません。話題は3つ並んでいますが、\textbf{読む人が判断すべきことが1つも書かれていない}からです。

\begin{quote}
\textbf{【定義:主張】}\ 主張とは、相手に受け入れてほしい、たった1つの文である。事実の列挙ではなく、判断である。承認・決定・同意という\textbf{相手の行動}に結びつく形をとる。
\end{quote}

「問い合わせが3件増えた」は事実です。「増加の要因を1名で調査する。よって予算を5万円使う」が主張です。報告書の最後には、必ずこの形の1文を置きます。

\subsection{論点と話題の違い}

よくある誤解は、話題が決まっていれば議論も進んでいると思うことです。両者は別物です。

\begin{itemize}
\item \textbf{話題}: 何の話をするか。案件名・会議体・資料のタイトル。
\item \textbf{論点}: 合意すべき判断が1つ決まっているか。Yes か No か、A か B か。
\end{itemize}

\begin{quote}
\textbf{【例】}\ 定例会議のアジェンダに「リリース日」と書かれたら、これは話題です。「7月1日で確定してよいか」まで書けて、初めて論点になります。前者の会議では、全員が別々の論点を頭に持って話します。
\end{quote}

\begin{quote}
\textbf{【注意】}\ 「その件について」「状況の共有です」で始まる議題は、たいてい論点が空です。空の論点は、その場の気分と発言力で埋まります。
\end{quote}

\subsection{絞る手順は3歩}

\begin{enumerate}
\item \textbf{相手に求める行動を1つ書く}: 承認する、了承する、選ぶ、期限を守る。
\item \textbf{その行動の理由を1文にする}: 「なぜ今、なぜこの案か」。
\item \textbf{残りはすべて根拠と反論への答えにする}: 主張を補強しない材料は、本文から外す。
\end{enumerate}

主張が2つ並ぶと、相手は弱いほうだけ潻して終われます。2つ欲しければ、2度書く前に\textbf{優先順位}を先に書きます。

\subsection{この本の地図}

この本は9章です。前半の第2〜4章で主張の部品を組み、中盤の第5〜6章で相手の攻撃から守り、後半の第7〜8章で実際の文と場面に載せます。最終章は提出前の点検表です。順に読めば、報告・提案・メール・会議の型が1つずつ手に入ります。

\section{事実と意見を分ける — 観測できたものと、解釈したものを切る}

議論が壊れる瞬間の多くは、事実と解釈が1文の中に入り込んだときです。分かれていない文は、根拠が崩れると解釈まで一緒に落ちます。

\begin{quote}
\textbf{【定義:事実と意見】}\ \textbf{事実}は、誰が観測しても同じ結果になるもの。記録・ログ・数値・発言の引用に残る。\textbf{意見(解釈)}は、その事実から自分が導いた意味づけである。
\end{quote}

\begin{quote}
\textbf{【例】}\ 「3月の問い合わせが3件増えた」は事実です。「顧客が不満を持っている」は解釈です。増えた原因が問い合わせ以外にある可能性、たとえば3月だけ案内メールを増やした事実を無視しています。
\end{quote}

事実は壊れません。壊れるのは解釈です。だから、両方を1文に混ぜた途端、\textbf{解釈に反論されると、事実の側の信用まで疑われる}。

\subsection{分ける言い方}

言い方には決まり文句があります。そのまま使います。

\begin{itemize}
\item 他人の発言・報告は「A社の担当者は〜と述べた」「〜という報告があった」。
\item 自分の判断は「私は〜と判断する」「〜と考える。理由は次節のとおり」。
\item 推定は「〜と推定される。ただし根拠は○○に限られる」。
\item 不確かな数値は「〜と見込まれる(出所:△△)」で、断定を避ける。
\end{itemize}

\begin{quote}
\textbf{【例】}\ 報告書の1節。「4月1日に障害が発生し、復旧は同日14時20分(監視ログより)。私は原因はメモリ不足と判断する。ただし根本原因の調査は未了であり、確定は5月上旬の見込みだ。」
\end{quote}

事実と判断が行で切られています。読者は事実だけを信じる選択も、判断だけを保留する選択もできます。

\begin{quote}
\textbf{【注意】}\ 事実のふりをした解釈は、見抜かれた瞬間、その報告全体の信用を失わせます。「社内で共通認識になっている」「多くの人がそう感じている」は、たいてい意見です。誰が何件、いつ、どこで言ったかが書ければ事実です。
\end{quote}

\begin{quote}
\textbf{【演習】}\ 次の4文を、事実・意見に分けよ。(1) 営業担当が先方の担当者から納期を3週間延長するよう求められた。(2) 顧客は当社に不満を持っている。(3) 昨月の解約理由の上位は「機能不足」で23件。(4) このままだと他社に流れる。
\end{quote}

\textbf{解答の指針:}\ (1) 事実。誰の何と言ったか、という観測記録だからだ。(2) 意見。「不満」という解釈に飛躍しており、解約理由のデータと一致するか確認が要る。(3) 事実。件数と分類が出所を伴っている。(4) 意見。未来の予測であり、「どの顧客が何件流れるか」の根拠が要る。答えを分ける基準は「記録に残るか」の1点である。

\section{根拠の3種 — データ・事例・権威}

主張を支える材料は、実務で使う限り3種類に分かれます。それぞれに得意な場面と、決まって崩される箇所があります。

\subsection{データ}

数値は、再現できるという最大の強みを持ちます。一方、潰され方は決まっています。\textbf{母数・期間・集計の定義}の3点です。

\begin{quote}
\textbf{【例】}\ 「問い合わせが20%増えた」と書けば「分母は?」「前月比か前年比か」「20%は何件か」と聞かれます。あらかじめ「前月比・4月は87件、3月は72件(社内トラッキングより)」まで書けば、この攻撃は閉じます。
\end{quote}

対策は1つ、\textbf{数字は単体で出さない}こと。母数、比較対象、期間、出所を数字の隣に置きます。

\subsection{事例}

具体例は、読者に場面を想像させます。説得力の体感がいちばん早い材料です。潰され方は「それは例外ではないか」です。

\begin{quote}
\textbf{【例】}\ 「A社で成功しています」への反論はすぐ出ます。「A社は導入担当に専任がいた。うちにはいない」。事例の条件と自分の条件がずれていると、効果は移りません。
\end{quote}

対策は、事例を出すとき\textbf{条件も出す}こと。担当体制、期間、予算。加えて、その事例が代表例か例外かを1行で書きます。

\subsection{権威}

第三者の判断は、相手の判断を速くします。潰され方は2点、\textbf{専門範囲}と\textbf{利害}です。

\begin{quote}
\textbf{【例】}\ 「著名な研究者が認めている」への反論は「その研究者はこの分野の専門か」「この製品を推薦する立場ではないか」です。権威の名前だけが独り歩きすると、指摘された瞬間に根拠が消えます。
\end{quote}

対策は、権威を出すとき「誰が」「何の分野で」「何と言ったか」「利害はあるか」の4点をセットにします。社内の権威なら「部長」だけでなく、部長の判断の理由を1行添えます。

\subsection{根拠の足りない主張の見つけ方}

見分け方は機械的です。\textbf{主張と根拠の間に「だから」を挟めるか}を確かめます。

\begin{quote}
\textbf{【例】}\ 「競合もやっている。市場も大きい。だからやるべきだ」。この文に「だから」は入りません。「競合が成功した理由が自社にも当てはまるから、やるべきだ」になって初めてつながります。並んでいるだけの材料は、根拠ではありません。
\end{quote}

自分の文に「だから」が挟めない箇所、相手の文に挟めない箇所。どちらも第1章の主張を支えていない箇所です。足りない主張は、\textbf{削るか、足すか}の二択で処理します。

\section{演繹と帰納を仕事の言葉で — 全部通る型と、例から拾う型}

推論のやり方は2つだけです。仕事の場面では、記号ではなく次の言い方で区別します。

\subsection{演繹型: すべてのAはB}

先に全般の規定、そこに個別の事例を当てはめる型です。

\begin{quote}
\textbf{【例】}\ 「顧客データは社外持ち出し禁止(規程)」「この作業で顧客データを私用端末にコピーした」「よってこの作業は規程違反」。前提が正しければ、結論は崩れません。
\end{quote}

強みは確実さです。潰され方は1か所、\textbf{前提}だけです。相手は「その規程は今の案件に適用されるのか」「例外規定はないのか」と攻めます。契約、規程、SLA、仕様。全称の規定がある場面で演繹型を使います。

\begin{quote}
\textbf{【注意】}\ 「すべて」「必ず」と書いたら、その1語全体に責任を持つことになります。例外を1つ知っている相手の前で、その1語は一撃で死にます。
\end{quote}

\subsection{帰納型: 例は3つある}

個別の観測から、傾向として一般化する型です。

\begin{quote}
\textbf{【例】}\ 「直近3件の障害はすべて同じバッチの順序起動が原因だった。よってこのバッチの順序を改めるべきだ」。例が3つあります。
\end{quote}

強みは、まだ規定のない新しい問題に使えることです。潰され方は「例外はないか」「たまたまではないか」です。だから帰納型では、\textbf{例の数と、例外の扱い}を先に書きます。「3件中3件で同一」「例外は未確認」と書けば、聞かれる前に答えていることになります。

どちらが有効かは、相手に何を聞きたいかで決まります。「例外は?」と聞いてほしい場面は演繹型で前提を揃え、「方向としてはこうだ」と切り開きたい場面は帰納型で例を積みます。

\subsection{よくある混同}

「例が3件ある」から「必ず起きる」への飛躍は、仕事の文で最も多い誤りです。3件は傾向の材料であって、保証ではありません。保証が要る場面(障害の再発防止、契約の遵守)では、例ではなく規程や仕様に足場を移します。

\section{反論の先回り — 棚卸しと、本文に埋める}

反論は、会議のその場で来るより、読んでいる最中に頭の中で来ます。先回りとは、その\textbf{予想}を書く前と書いた直後の2か所に配置する作業です。

\subsection{棚卸し: 相手が突く箇所を先に書く}

主張を書き終えたら、紙を1枚、相手の立場で並べます。使う問いは決まっています。

\begin{itemize}
\item \textbf{真偽}: この根拠は本当か。出所はどこか。
\item \textbf{十分性}: それだけで足りるか。他に要因はないか。
\item \textbf{今か}: なぜ今か。1か月後でも間に合うのではないか。
\item \textbf{コスト}: 時間とお金と人員は。誰がやるのか。
\item \textbf{失敗}: 失敗したらどうなるか。戻せるか。
\end{itemize}

5つの問いに全部答える必要はありません。\textbf{自分の主張の中で最も脆い2つ}だけを選び、そこに答えを書きます。

\subsection{回答は本文に埋め込む}

回答を最後の「FAQ」にまとめると、相手はFAQだけ読んで反論します。埋め込み型のほうが強い。

\begin{quote}
\textbf{【例】}\ 「導入コストは120万円で、単月の回収は8か月と見込む。初期費用が重いという指摘には、2部門への段階導入で初月40万円に分割でき、残りは効果が出た第2月以降に回すと応える。」
\end{quote}

反論への回答が、主張の直後に入っています。読者は次の段落に進む前に、その反論を手放します。

口頭では同じ型を最初の1文で使います。「おっしゃる通り、初期費用は懸念です。その点は段階導入で応えます」。相手を一度認めてから戻すと、会議は前進します。

\begin{quote}
\textbf{【注意】}\ 反論を全部消すのは狙い違いです。\textbf{回答できる反論だけ}先に置きます。答えられない反論は、主張の範囲を狭めて消すか、「未対応の懸点」として明示します。隠した反論は、いちばん先に掘り出されます。
\end{quote}

\begin{quote}
\textbf{【演習】}\ 「取引先への新しい請求方式への切り替え」を提案する文を1段落書き、次に想定される反論を2つ挙げ、それぞれへの回答を本文に埋め込む形に書き直せ。
\end{quote}

\textbf{解答の指針:}\ 反論候補は「先方の経理の作業量」「移行中の二重請求の取り扱い」「開始時期」などから探す。埋め込みの形は「切り替えは10月から。ただし二重請求が出た場合の精算ルールを同封し、作業増は月2時間までに抑える」のように、主張・懸念・回答を1段落に収める。反論が本文に見当たらない文は、棚卸しが足りない。

\section{よくある崩し方 — 転換・藁人形・数の暴力・権威の盾}

会議で主張が崩れる瞬間には、決まった型があります。名前を知っておくと、その場で戻せます。自分が使わないことも同じほど大事です。

\subsection{転換: 問いを別の問いにすり替える}

質問を、答えると都合のよい別の質問に入れ替える崩し方です。

\begin{quote}
\textbf{【例】}\ 問い「3月だけ件数が増えた要因は?」への答え「当社の集計は常に厳格です」。集計の厳格さは問われておらず、3月の例外要因が問われています。
\end{quote}

受け方は1つ、\textbf{問いをもう一度言ってから答えます}。「いま問われているのは集計の厳格さではなく、3月の例外要因です」。問いを取り直せば、転換は自然に切れます。

\subsection{藁人形: 弱く言い替えて叩く}

相手の主張を弱い表現に置き換えてから倒す崩し方です。

\begin{quote}
\textbf{【例】}\ 主張「予算を20%増やして人手を補う」への応答「つまり金が足りないと、昔から言っているだけですね」。予算の内訳と回収計画は論点に上がっていません。
\end{quote}

受け方は、\textbf{自分の主張を1文で再提示して戻す}ことです。「私の主張は、増収分の20%を人件費に回すことです。昔からの一般論とは別の計画を添えます」。言い換えさせないことが防御になります。逆に、相手の主張を勝手に短く要約してから反論してはいけません。

\subsection{数の暴力: 多さで押す}

数の多い側が正しいと考えさせる崩し方です。

\begin{quote}
\textbf{【例】}\ 「反対は20件、賛成は3件だから多数決は明らかだ」。母数が知られません。1人が何件も出している可能性、重複、社内アンケートの回収率5%という条件も伏せられています。
\end{quote}

受け方は、\textbf{母数・重複・条件を聞く}だけです。「全員何名中か」「1人何件までか」。2つの問いで、数字の見え方は変わります。自分の側の数字も同じ検査を受けます。

\subsection{権威の盾: 名前で議論を止める}

権威を、議論を遮る盾として使う崩し方です。第3章の権威は根拠として足を前に向かせますが、この権威は足を止めます。

\begin{quote}
\textbf{【例】}\ 「部長もその方向で話している」で議論が終わる場面。部長の判断の理由が共有されていなければ、部長の名前は根拠ではなく休止符です。
\end{quote}

受け方は、\textbf{権威の主張の内容と、その理由を分けて聞く}ことです。「部長の方針の理由は、コストか、納期か、どちらでしょうか」。理由が出てくれば議論は再開し、出なければ「理由の共有をお願いします」と記録に残します。

\begin{quote}
\textbf{【演習】}\ 会議で「反対が20件も出ている。賛成3件の案を続ける理由はどこにある?」と指摘された。まず確認すべきことを2つ挙げ、その上で自分の提案の範囲を限定して応じる台詞を書け。
\end{quote}

\textbf{解答の指針:}\ まず「全員中何名か」と「1人何件出しているか」の2つを確認する。その上で「反対20件のうち、重複と未回答を除くと○件。この案はA部門に限定した試行で、範囲外の反対は移行判断で考慮する」のように、母数の確認と主張の範囲限定を1台詞に収める。数の暴力への応答は、数字の否定ではなく条件の追加で成立する。

\section{構成の型 — PREP、箇条書きの順、見出しと主張}

部品がそろっても、並びが悪いと主張は届きません。実務で一番安定する並びがPREPです。

\subsection{PREP: 結論・理由・例・再結論}

\begin{itemize}
\item \textbf{Point}: 主張を1文。最初の一文に置く。
\item \textbf{Reason}: 理由を1つか2つ。「なぜなら」でつなぐ。
\item \textbf{Example}: 具体例・数値・場面を1つ。
\item \textbf{Point}: 主張を言い直し、相手に求める行動を1つ添える。
\end{itemize}

\begin{quote}
\textbf{【例】}\ 「問い合わせ対応の一次回答を自動化すべきだ(結論)。一次回答の7割は定型で、担当者の確認を要さないからだ(理由)。実際、4月に自動化した返信テンプレートは、適用の82%で追加対応なしで終わった(例)。よって定型30種に自動化を適用し、来月の運用会議で了承を得たい(再結論と行動)」。
\end{quote}

結論を先に置くのは、\textbf{相手が判断する係}だからです。理由から始めた文は、相手が結論を待たずに反論の準備を始め、聞かれていない質問に答えてしまうことになります。

\subsection{箇条書きの論理的並び順}

箇条書きには順序があり、\textbf{順序には名前が付きます}。

\begin{itemize}
\item \textbf{時系列}: 事実が起きた順。障害報告や経緯説明に使う。
\item \textbf{因果}: 原因から結果へ。分析や要因究明に使う。
\item \textbf{重要度}: 判断に効くものを先頭へ。提案や選択肢比較に使う。
\end{itemize}

重要度順に並べるときは、先頭の項目こそが主張に最も近い項目です。逆に、重要度不明のまま並べると、読者は「どれが本命か」を測るために全部を精読します。その負担が、差し戻しの原因になります。

\subsection{見出しと主張の対応}

見出しは、その節に書かれることへの約束です。約束を判断文にできれば、\textbf{目次だけで主張が読めます}。

\begin{quote}
\textbf{【例】}\ 「背景」「現状」「詳細」は情報の箱です。「A案の採用理由」「7月開始が妥当な3つの根拠」は判断です。後者の目次を読むだけで、承認するかどうかの材料が手に入ります。
\end{quote}

\begin{quote}
\textbf{【注意】}\ PREPの最後のPointは、要約ではありません。\textbf{決定の依頼}です。「以上です」で終えず、「どの案に了承いただくか」まで書きます。行動が1つもない再結論は、読み終えた相手を置き去りにします。
\end{quote}

\begin{quote}
\textbf{【演習】}\ 次の4つの文をPREPの順に並べ、Pointの文を「相手の行動」が入る形に書き直せ。「導入は9月が最適だ」「経理の作業は月4時間に収まる」「費用は月8万円で、対応件数から割ると1件240円だ」「この方式を採用したい」。
\end{quote}

\textbf{解答の指針:}\ 順は「この方式を採用したい」(Point)→「費用は月8万円、1件240円」(Reason)→「経理の作業は月4時間に収まる」(Example相当の具体)→「導入は9月が最適だ」(再Point)。最後のPointには「9月開始で了承いただきたい」という行動が要る。理由と例の区別が曖昧でも、先頭と末尾に判断が来る形ならPREPの骨格になる。

\section{会議とメールでの運用 — 質問・同意・判断の止め方}

同じ論理が、場面で使えなければ意味がありません。質問への答え方、同意の取り方、判断を止める問いの3つを押さえます。

\subsection{質問への答え方}

質問には種類があり、種類ごとに答え方が違います。

\begin{itemize}
\item \textbf{事実の質問}: 即答する。出所も添える。「72件です。トラッキングの4月分より」。
\item \textbf{判断の質問}: 条件つきで答える。「予算次第で10月も可能ですが、9月なら回収計画が1か月前倒しになります」。
\item \textbf{意見の質問}: 自分の立場と根拠を分けて出す。「私は賛成です(判断)。理由は件数の半減が観測されていること(根拠)」。
\end{itemize}

答えられない質問には、\textbf{答えられないという答え}が正当です。「分かりません。金曜までに調査して、同じこの場にお持ちします」。お茶を濁すより、期限付きの持ち帰りが信用します。

\subsection{同意の取り方}

一括での同意は、後で崩れます。合意の単位を小さく切ります。

\begin{enumerate}
\item \textbf{論点の合意}: 今日決めるのは何かを先に言葉にする。
\item \textbf{主張の合意}: その決めるべき1つに対して、自分の案を置く。
\item \textbf{根拠の合意}: 理由と例に、納得いかない点がないかを確認する。
\end{enumerate}

\begin{quote}
\textbf{【例】}\ 「ここまでは、障害の原因はバッチ順序で合意ですね。では対応は今週か来週か、の2案で考えます」。ここまで区切って発言すると、後から「原因はまだ分からない」と言われても、合意済みの範囲が残ります。
\end{quote}

\begin{quote}
\textbf{【注意】}\ 「合意したつもりの8割は、書いていない合意です」。決めたことは、その場で1行、議事に残します。書かれていない合意は、関係者が替わった瞬間に消えます。
\end{quote}

\subsection{判断を止める問い}

会議が空転するときは、問いが複数になっているのが原因です。問いを1つに戻します。

\begin{quote}
\textbf{【例】}\ 「今日決めたいのは、予算を取るか、取らずに来期に回すかのどちらかです。手順の細かいことは、決定後に詰めます」
\end{quote}

決められない日は、決めないことも明示します。「これは今日は決めません。資料が揃う5月の会議に持ち帰ります」。持ち帰りを宣言しない会議は、終了時も同じ問いのままです。

メールでの運用も同じです。\textbf{1通1主張}にします。複数の相手への複数の依頼は、通ごとに分けます。返信を求めるなら期限を書き、それまでに返事が無いときに何をするかも添えます。長いスレッドは、最新の返信に結論を再掲します。過去ログを読ませる義務を、読者に負わせないためです。

\section{崩されないチェックリスト — 提出・発言前の点検12項目}

最後は、点検表です。提出前、発言前、送信前に上から順に確認します。1つでも引っかかったら、そこだけ直してから出します。

\begin{enumerate}
\item 主張を1文で書けるか。
\item 相手に求める行動が1つに決まっているか。
\item 話題ではなく論点を書いたか(Yes か No かが分かる形か)。
\item 事実と意見が、文または行で分かれているか。
\item データに母数・期間・出所が付いているか。
\item 事例は、代表例か例外かが明示されているか。
\item 権威は、専門範囲と利害の確認を経ているか。
\item 「すべて」「必ず」の責任を取れるか。例外を知られていないか。
\item 想定反論2つ以上に、本文で答えているか。
\item 転換・藁人形・数の暴力・権威の盾を、自分が使っていないか。
\item PREPになっており、目次だけで主張が読めるか。
\item 合意の単位を区切り、未決事項と持ち帰りを明記しているか。
\end{enumerate}

12項目は、この本の前半から順に対応しています。1〜3は第1章、4は第2章、5〜7は第3章、8は第4章、9は第5章、10は第6章、11は第7章、12は第8章です。どこで引っかかったかは、直す場所を教えてくれます。

\begin{quote}
\textbf{【総合演習】}\ 次の提案メールを12項目で採点し、まず3つの不合格を挙げ、本文を修正せよ。「新ツールの導入を検討しました。便利そうなので導入したいです。他社も導入しています。また、コストも安いです。弊社の業務にも合うと思います。ぜひご検討ください。」
\end{quote}

\textbf{解答の指針:}\ まず挙げるのは、(1)主張と行動の不定(「検討」か「導入の承認」かが不明、項目1・2)、(2)根拠の欠如(「便利そう」「安い」に数値と出所がなく、「他社も」は条件の合う事例か不明、項目5・6)、(3)反論への回答ゼロ(移行コスト、既存ツールとの重複に触れていない、項目9)。修正は「月額と移行期間を明記したうえで、まず1部門で3か月の試用の承認をいただきたい」を冒頭に置き、後半を理由と例に置き換える形で合格圏に入る。

論理の組み立ては、相手を黙らせる技術ではありません。\textbf{相手が判断しやすくする装置}です。主張を1つに絞り、根拠を分けて、反論を先に受ける。この3つが身につけば、報告も提案も、会議もメールも、同じ形で組み上がります。

\end{document}
